\input{preamble.tex}

\title{Лабораторная работа № 7. \\ Учёт физических параметров сети}
\author{Сущенко Алина Николаевна \\ НПИбд-01-23}

\institute{Российский университет дружбы народов имени Патриса Лумумбы}
\date{2026}

\begin{document}

\frame{\titlepage}

\begin{frame}
\frametitle{Цель работы}
\begin{itemize}
    \item Приобретение практических навыков по работе с физической рабочей областью Packet Tracer.
\end{itemize}
\end{frame}

\begin{frame}
\frametitle{Настройка физической рабочей области}
В физической рабочей области Packet Tracer присвоили название городу - \texttt{ansusenko}.
    \centering
    \includegraphics[width=\textwidth]{../images/01.png}
    \captionof{figure}{Присвоение названия городу ansusenko.}
\end{frame}

\begin{frame}
\frametitle{Настройка физической рабочей области}
Щёлкнув на изображении города, перешли на уровень зданий. Присвоили зданию название \texttt{ansusenko}. Добавили здание для территории \texttt{ansusenko-p}.
    \centering
    \includegraphics[width=\textwidth]{../images/02.png}
    \captionof{figure}{Здания ansusenko и ansusenko-p.}
\end{frame}

\begin{frame}
\frametitle{Настройка физической рабочей области}
Щёлкнув на изображении здания ansusenko, переместили изображение, обозначающее серверное помещение (Main Wiring Closet), в него.
    \centering
    \includegraphics[width=\textwidth]{../images/03.png}
    \captionof{figure}{Размещение серверного помещения в здании ansusenko.}
\end{frame}

\begin{frame}
\frametitle{Настройка физической рабочей области}
    \centering
    \includegraphics[width=\textwidth]{../images/04.png}
    \captionof{figure}{Размещение сетевого оборудования в серверной стойке.}
\end{frame}

\begin{frame}
\frametitle{Настройка физической рабочей области}
На территории \texttt{ansusenko-p} разместили коммутатор \texttt{msk-ansusenko-sw-1} и оконечные устройства \texttt{other-ansusenko-1}, \texttt{dk-ansusenko-1}.
    \centering
    \includegraphics[width=\textwidth]{../images/05.png}
    \captionof{figure}{Размещение оборудования на территории ansusenko-p.}
\end{frame}

\begin{frame}
\frametitle{Настройка IP-адресации}
    \centering
    \includegraphics[width=\textwidth]{../images/06.png}
    \captionof{figure}{Пинг на коммутаторе msk-ansusenko-sw-01.}
\end{frame}

\begin{frame}
\frametitle{Настройка IP-адресации}
    \centering
    \includegraphics[width=\textwidth]{../images/07.png}
    \captionof{figure}{Неудачный пинг.}
\end{frame}

\begin{frame}
\frametitle{Проверка работоспособности соединения}
    \centering
    \includegraphics[width=\textwidth]{../images/08.png}
    \captionof{figure}{Неудачный пинг после изменения расстояния.}
\end{frame}

\begin{frame}
\frametitle{Добавление повторителей}
    \centering
    \includegraphics[width=\textwidth]{../images/09.png}
    \captionof{figure}{Вид повторителя с установленными модулями.}
\end{frame}

\begin{frame}
\frametitle{Схема сети с повторителями}
    \centering
    \includegraphics[width=\textwidth]{../images/10.png}
    \captionof{figure}{Новый вид схемы (отрезок с повторителями).}
\end{frame}

\begin{frame}
\frametitle{Изменение расстояния между территориями}
    \centering
    \includegraphics[width=\textwidth]{../images/11.png}
    \captionof{figure}{Расстояние между территориями 991.37 м, соединены оптоволокном.}
\end{frame}

\begin{frame}
\frametitle{Проверка работоспособности соединения после добавления повторителей}
    \centering
    \includegraphics[width=\textwidth]{../images/12.png}
    \captionof{figure}{Ping от ansusenko к ansusenko-p.}
\end{frame}

\begin{frame}
\frametitle{Проверка работоспособности соединения после добавления повторителей}
    \centering
    \includegraphics[width=\textwidth]{../images/13.png}
    \captionof{figure}{Ping от ansusenko-p к ansusenko.}
\end{frame}

\begin{frame}
\frametitle{Контрольные вопросы}

\begin{block}{Вопрос 1. Перечислите возможные среды передачи данных. На какие характеристики среды передачи данных следует обращать внимание при планировании сети?}
\textbf{Ответ:} \\
\textbf{Среды передачи данных:}
\begin{itemize}
    \item \textbf{Проводные (медные):} витая пара (UTP, FTP, STP), коаксиальный кабель.
    \item \textbf{Оптические:} одномодовое оптоволокно (SMF), многомодовое оптоволокно (MMF).
    \item \textbf{Беспроводные:} радиоканал (Wi-Fi, Bluetooth, LTE), инфракрасный диапазон.
\end{itemize}
\textbf{Характеристики для планирования:}
\begin{itemize}
    \item Максимальная длина сегмента (затухание сигнала)
    \item Пропускная способность (скорость передачи данных)
    \item Помехозащищённость (электромагнитные помехи, перекрёстные наводки)
    \item Стоимость прокладки и обслуживания
    \item Сложность монтажа и масштабирования
\end{itemize}
\end{block}

\end{frame}

\begin{frame}
\frametitle{Контрольные вопросы}

\begin{block}{Вопрос 2. Перечислите категории витой пары. Чем они отличаются? Какая категория в каких условиях может применяться?}
\textbf{Ответ:} \\
\textbf{Категории витой пары:}
\begin{itemize}
    \item \textbf{Category 5e (Cat5e):} частота до 100 МГц, скорость до 1 Гбит/с (1000BASE-T). Применяется в офисных сетях, домашних сетях.
    \item \textbf{Category 6 (Cat6):} частота до 250 МГц, скорость до 10 Гбит/с на расстоянии до 55 м. Применяется в корпоративных сетях, дата-центрах.
    \item \textbf{Category 6a (Cat6a):} частота до 500 МГц, скорость до 10 Гбит/с на расстоянии до 100 м. Экранированная версия, устойчива к помехам.
    \item \textbf{Category 7 (Cat7):} частота до 600 МГц, полностью экранированная (S/FTP), используется в высоконагруженных сетях.
    \item \textbf{Category 8 (Cat8):} частота до 2000 МГц, скорость до 25-40 Гбит/с, применяется в дата-центрах (расстояние до 30 м).
\end{itemize}
\textbf{Отличия:} рабочая частота, максимальная скорость передачи, наличие экранирования, максимальная длина сегмента.
\end{block}

\end{frame}

\begin{frame}
\frametitle{Контрольные вопросы}

\begin{block}{Вопрос 3. В чем отличие одномодового и многомодового оптоволокна? Какой тип кабеля в каких условиях может применяться?}
\textbf{Ответ:} \\
\textbf{Одномодовое оптоволокно (SMF):}
\begin{itemize}
    \item Диаметр сердцевины: 8-10 мкм
    \item Распространяется только одна мода света
    \item Дальность передачи: до 100 км и более
    \item Применение: магистральные линии связи, соединение городов, дата-центры на больших расстояниях
\end{itemize}
\textbf{Многомодовое оптоволокно (MMF):}
\begin{itemize}
    \item Диаметр сердцевины: 50-62.5 мкм
    \item Распространяется множество мод света
    \item Дальность передачи: до 550 м (для 10 Гбит/с), до 2 км (для 1 Гбит/с)
    \item Применение: локальные сети внутри зданий, кампусные сети, соединение этажей
\end{itemize}
\textbf{Основное отличие:} диаметр сердцевины, дальность передачи, стоимость (SMF дороже в оборудовании, MMF — в кабеле).
\end{block}

\end{frame}

\begin{frame}
\frametitle{Контрольные вопросы}

\begin{block}{Вопрос 4. Какие разъёмы встречаются на патчах оптоволокна? Чем они отличаются?}
\textbf{Ответ:} \\
\textbf{Основные типы оптических разъёмов:}
\begin{itemize}
    \item \textbf{SC (Subscriber Connector):} квадратный форм-фактор, фиксация "push-pull". Используется в Ethernet-сетях, простота подключения.
    \item \textbf{LC (Lucent Connector):} компактный разъём (половина SC), фиксация защёлкой. Самый распространённый в современных SFP-модулях.
    \item \textbf{ST (Straight Tip):} круглый разъём с байонетным креплением (как BNC). Используется в старых сетях.
    \item \textbf{FC (Ferrule Connector):} круглый разъём с винтовой фиксацией. Применяется в условиях вибрации (промышленность, телекоммуникации).
    \item \textbf{MTP/MPO (Multi-fiber Push On):} многожильный разъём (до 24 волокон). Используется в дата-центрах для высокой плотности подключений.
\end{itemize}
\textbf{Отличия:} форма разъёма, способ фиксации, плотность портов, область применения.
\end{block}

\end{frame}

\begin{frame}
\frametitle{Выводы}
\begin{itemize}
    \item В результате выполнения лабораторной работы приобрели практические навыки по работе с физической рабочей областью Packet Tracer.
\end{itemize}
\end{frame}

\end{document}
